\section{References}

Sources are grouped by the role they play. Only sources actually read for this edition are
listed, and each entry states what was read.

\subsection{Promulgating acts}

\begin{itemize}[leftmargin=1.3em, itemsep=0.3em]
\item John XXIII, apostolic letter given \latin{motu proprio}, \latin{Novum rubricarum Breviarii
et Missalis Romani corpus approbatur} (\latin{Rubricarum instructum}), 25 July 1960. \latin{Acta
Apostolicae Sedis} 52 (1960): 593--595. The complete text was read; the preamble and the six
numbered dispositions are used throughout Section~1. The same text is reprinted in the 1962
typical edition of the Missal, where it likewise carries the date 25 July 1960.
\item Sacred Congregation of Rites, general decree \latin{Quo novus rubricarum Breviarii ac
Missalis Romani codex promulgatur} (\latin{Novum rubricarum}), 26 July 1960. \latin{Acta
Apostolicae Sedis} 52 (1960): 596. The complete text was read; its provision annexing the
\latin{Variationes} is used in Sections~1 and~6.
\item Sacred Congregation of Rites, decree declaring the new Vatican edition of the Roman Missal
typical, 23 June 1962, subscribed by A. M. Cardinal Larraona, Prefect, and Enrico Dante,
Secretary. Printed at the head of the 1962 typical edition. The complete text was read; its
recitation of the date of \latin{Rubricarum instructum} is treated in Section~\ref{sec:collation}.
\end{itemize}

\subsection{The controlling norms}

\begin{itemize}[leftmargin=1.3em, itemsep=0.3em]
\item \latin{Rubricae Breviarii et Missalis Romani}, 26 July 1960. \latin{Acta Apostolicae
Sedis} 52 (1960): 597--740. \latin{Pars prima, Rubricae generales}, \rnn{1--137}, was read in
full at pages 597--621 and is the principal source of Sections~2 through~5 and~7 through~8.
\latin{Pars tertia, Rubricae generales Missalis Romani}, was read at the chapters bearing on the
calendar and the orations: \rnn{269--279} on the calendar to be used in celebrating Mass;
\rnn{298--305} on the Mass on Sundays, ferias, and feasts; \rnn{341--347} on votive Masses of the
II class and the Rogation Mass; and \rnn{433--465} on the orations. Printed pages 598 and 599
were additionally read as rendered page images at 400 dots per inch.
\item The same code as reprinted in \latin{Missale Romanum}, editio typica (Vatican City: Typis
Polyglottis Vaticanis, 1962), pages XIII--XXXVI, carrying \latin{pars prima} and \latin{pars
tertia} only, as \rn{5} of \latin{Rubricarum instructum} requires. The whole of \latin{pars
prima} was read, and the pages carrying \rnn{1--19} and the table at \rn{91} were read as
rendered page images at 300 dots per inch.
\item \latin{Variationes in Breviario et Missali Romano ad normam novi codicis rubricarum},
annexed to the promulgating decree. \latin{Acta Apostolicae Sedis} 52 (1960): 706 and following.
\latin{Caput I, Variationes in Calendario}, \rnn{1--12}, was read in full and supplies the
itemised calendar changes in Section~\ref{sec:sanctoral}; \rn{13} of \latin{Caput II} supplies
the abolition of the Suffrage of All Saints and the commemoration of the Cross in
Section~\ref{sec:commemorations}.
\end{itemize}

\subsection{The typical edition and its calendar}

\begin{itemize}[leftmargin=1.3em, itemsep=0.3em]
\item \latin{Missale Romanum ex decreto Sacrosancti Concilii Tridentini restitutum, Summorum
Pontificum cura recognitum}, editio typica (Vatican City: Typis Polyglottis Vaticanis, 1962).
Read in the facsimile published by the Church Music Association of America, whose exact bytes
are registered in the repository's source library and were re-downloaded and hashed
byte-identical during the preparation of this edition. The following were read:
\begin{itemize}[leftmargin=1.2em, itemsep=0.1em, topsep=0.15em]
\item the front-matter decrees and papal bulls;
\item the note \latin{De anno et eius partibus}, including the Ember Day rule and the account of
the Gregorian correction;
\item the \latin{Tabella temporaria festorum mobilium} and the \latin{Tabula paschalis antiqua
reformata}, read as rendered page images at 200 and 400 dots per inch;
\item the universal \latin{Calendarium} at printed pages XLV--LIII, all nine pages read as
rendered page images at 260 dots per inch;
\item the formularies of the First Sunday after the Epiphany, the Holy Family, and the
Commemoration of the Baptism of the Lord, with their occurrence rubrics;
\item the Ember Wednesday of Pentecost and of September;
\item the Twenty-third Sunday after Pentecost with the resumption rubric printed at its end;
the Third and Fourth Sundays left over after the Epiphany; and the Twenty-fourth and last Sunday
after Pentecost, whose Epistle is Colossians 1:9--14 and Gospel Matthew 24:15--35.
\end{itemize}
\item Pius V, bull \latin{Quo primum tempore}, 14 July 1570; Clement VIII, brief of 7 July 1604;
and Urban VIII, brief of 2 September 1634 — all read in the front matter of the same 1962
edition, and used only for the dated timeline.
\end{itemize}

\subsection{Comparative witness, printing year unestablished}

\begin{itemize}[leftmargin=1.3em, itemsep=0.3em]
\item \latin{Missale Romanum \dots\ a Pio X reformatum, Benedicti XV auctoritate vulgatum},
\latin{editio IV iuxta typicam Vaticanam} (New York: Benziger Brothers, Inc., copyright 1942 and
1944), with the approval leaf of Francis Joseph, Archbishop of New York, dated 8 December 1942
and reciting conformity to the decree of the Sacred Congregation of Rites of 9 January 1942, and
with the Congregation's decree of 25 July 1920 declaring the Vatican edition typical. Read in an
Internet Archive optical text derivative whose exact bytes are registered in the repository's
source library; the item's ``1947'' is uploader-supplied and unconfirmed. It carries the
pre-1955 general rubrics with the double, semidouble, and simple grades. Four passages were
used: the resumption rubric printed at the end of the Twenty-third Sunday after Pentecost; the
note fixing the Ember Days by the third Sunday of Advent, the first Sunday of Lent, Pentecost
Sunday, and the feast of the Exaltation of the Holy Cross; \latin{Additiones et Variationes},
\latin{De occurrentia et de translatione festorum}, \rn{6}, on resuming an impeded Sunday's Mass
within the following week; and the ceiling of seven orations, of odd number, in the same
\latin{Additiones et Variationes}. It is a Missal and carries no Breviary rubrics and no tables
of occurrence or concurrence. The witness controls no reading of the 1962 books and no dating of
any change; see Section~\ref{sec:witness}.
\item \latin{Missale Romanum \dots\ editio iuxta typicam} (New York: Benziger Brothers, Inc.,
1962), whose approval leaf recites that the Sacred Congregation of Rites declared the Benziger
edition \latin{iuxta typicam} on 21 October 1961. Read in an Internet Archive optical text
derivative already registered in the repository's source library. One passage was used: the
front-matter note fixing the Ember Days, which reads \latin{post dominicam III septembris} and
so corroborates the 1962 Vatican wording from an independently produced book conformed to the
same code.
\end{itemize}

\subsection{Computation}

\begin{itemize}[leftmargin=1.3em, itemsep=0.3em]
\item The paschal dates, the derived movable anchors, the counts $E$ and $P$, and the surplus $S$
were computed for 1900--2100 from the Gregorian ecclesiastical rule and compared with the
\latin{Tabella temporaria} for the fifty-two years 1960--2011 printed on its first page. The
method, the agreement, and the single divergence are stated in Sections~\ref{sec:movable}
and~\ref{sec:collation}. No source outside the Missal is relied on for any dated result.
\item The repository's own \latin{Roman Calendar Computation} guidance owns the shared arithmetic
used across its liturgy collections. This edition states and applies those rules and does not
establish a divergent one; the open points it records — the absence of an official statement of
the computus, and the unlegislated shortfall case — are the same points recorded there.
\end{itemize}
